VERA adapts a frozen vision-language model by rewriting the questions it answers, and scores only the answer text. We asked whether InternVL2-8B, adapted this way to the VISILAB firearm-action dataset, decides a clip's category from the firearm, and whether telling VERA's optimizer where the model looked would make it rely on the firearm more. Two pre-registered measures were taken on all 398 videos with the untrained questions: the decision's attention to the annotated firearm relative to chance (lift), and the drop in the true category's probability when the firearm is greyed out rather than an equal patch elsewhere (reliance). The attention is specific to the firearm: lift exceeds that of another video's boxes by \val{gate_V1_estimate} [\val{gate_V1_low}, \val{gate_V1_high}]. Correct answers depend on the firearm entirely: every one of them changes when it is masked. But lift does not track reliance ($\rho=$ \val{gate_V2_estimate} [\val{gate_V2_low}, \val{gate_V2_high}]), so the pre-registered gate failed and the training experiment did not run. The model looks at handguns yet names only \val{R5_seed_Handgun_category_acc_pct}\,\% of the Handgun videos correctly, calling most of the rest No\_Gun: its failure is recognition, not where it looks, and attention feedback would not address it.
